\scanpage{282}
\mspage{198}
% archive sheet number: 281
Enfin b) $\Longrightarrow$ a) résulte du fait que $\mathrm{Ind}_{\unsure{\pi|\kappa|k}_0}(C)$ est stable par $\varinjlim$ \unsure{grdes|gdes} devant $\pi_0$ (4.4.15), et que $C$ est une sous-catégorie str. génératrice de $\mathrm{Ind}_{\pi_0}(C)$, formée d'objets $\pi_0$-accessibles.

\deleted{Pour vérifier que $C$ est strictement génératrice, en vertu de 4.8.2, $1^\circ$, il suffit, puisque $C$ est une sous-catégorie \uline{pleine} de $M=\mathrm{Ind}_{\pi_0}(C)$, de vérifier que pour \unsure{tt} \inserted{$F$ dans $M$, la} famille des flèches \deleted{de} $x\longrightarrow F$ de but $F$, de source dans $C$, est épim. stricte. Pour ceci, \deleted{il suffit d'ailleurs de le vérifier \unsure{dans}}}

\unsure{Il|il} nous faut expliciter le

\uline{Lemme} 4.8.4. Soit $C$ une \inserted{$\unsure{\mathcal U|U}$-}catégorie, \deleted{\unsure{égale|équi.}} \inserted{$\pi_0$ un cardinal infini} \deleted{\unsure{à une} petite} Alors $C$ est \unsure{strictement|stricte} génératrice dans \inserted{$M=$}$\mathrm{Ind}\, C$ et dans \inserted{$M=$}$\mathrm{Ind}_{\pi_0}C$.

\deleted{Il faut prouver, $F$ étant un objet de $M$, que}

\[ \deleted{\varinjlim_{C_{/F}}x\longrightarrow F} \]

\deleted{est un iso. On a $F=\varinjlim_{\deleted{\unsure{x}}{\unsure{I|\underline I}}}x_i$, $I$ \inserted{ens. ordonné} filtrant (\unsure{grd} devant $\pi_0$ dans le deuxième cas), les $x_i$ dans $C$, pour un foncteur $I\longrightarrow C\deleted{\unsure{_{/F}}}$ convenable. On a une factorisation \inserted{can.} \unsure{de}}
